% 04-main-result.tex
% Section 4 — Main result. Drafted G3.4.
% [ERIC-REVIEW] on mathematical phrasing + proof sketch accuracy.

\section{Main Result}
\label{sec:main}

We state the main theorem in its Lean 4 formulation and sketch the
structure of its proof. Detailed discussion of the hypotheses is deferred
to Section~\ref{sec:hypotheses}, and the formalisation choices are
described in Section~\ref{sec:lean}.

\subsection{Statement}

\begin{theorem}[Main theorem, conditional]\label{thm:main}
% [ERIC-REVIEW] This is the natural-language rendering of
% ProjetCollatz.no_nontrivial_cycle_phase59. Please verify the English
% matches the intended mathematical content.
Let \(\collatz : \Nat \to \Nat\) be the Collatz map, and call \emph{nontrivial
odd cycle} a pair \((n, k)\) with \(k \geq 2\) such that \((n, k)\)
satisfies the formal predicate \texttt{IsOddCycle} from the repository
(see Section~\ref{sec:preliminaries} and the Lean definition). Then, under
the three hypotheses described in Section~\ref{sec:hypotheses} ---
\textsf{BakerSeparation}, \textsf{BarinaVerification}, and
\textsf{DerivedLargeKBound} --- no such nontrivial odd cycle exists.
\end{theorem}

The formal Lean statement is :

\begin{lstlisting}[caption={\texttt{no\_nontrivial\_cycle\_phase59} (\texttt{Phase59ContinuedFractions.lean}:236).}]
theorem no_nontrivial_cycle_phase59
    (baker : BakerSeparation) (barina : BarinaVerification)
    (cf : DerivedLargeKBound)
    (n k : N) (hcyc : IsOddCycle n k) : False
\end{lstlisting}

The conclusion \texttt{False} from the assumption \texttt{IsOddCycle n k}
is equivalent to the non-existence of a witness : the theorem rules out
any pair \((n, k)\) that satisfies \texttt{IsOddCycle}.

\subsection{Proof structure}

The proof is a case split on the value of \(k\) relative to the threshold
\(1322\) introduced by Hypothesis~\textsf{DerivedLargeKBound}
(see Section~\ref{sec:hypotheses}). Writing the Lean proof term
schematically :

\begin{lstlisting}
theorem no_nontrivial_cycle_phase59 baker barina cf n k hcyc := by
  by_cases hk : k \leq 1322
  \cdot exact no_cycle_k_le_1322 baker barina n k hcyc hk
  \cdot push_neg at hk
    exact no_cycle_k_gt_1322 barina cf n k hcyc hk
\end{lstlisting}

The two branches are handled by the following sub-lemmas, both proved in
the repository with the same three-axiom chain as the main theorem
(\texttt{propext}, \texttt{Classical.choice}, \texttt{Quot.sound}) :

\begin{itemize}
  \item \textbf{\texttt{no\_cycle\_k\_le\_1322}} : the range \(2 \leq k
        \leq 1322\) is handled by combining a lower bound on the seed
        \(n\) of any hypothetical cycle (from the Product Bound argument
        built on Baker's theorem, see Section~\ref{sec:hypotheses}) with
        the Barina range \(n \geq 2^{71}\). Hypothesis
        \textsf{BarinaVerification} then yields a contradiction against
        any cycle with \(n < 2^{71}\), while the Product Bound shows
        \(n \geq 2^{71}\) under Baker's hypothesis ; the two bounds are
        incompatible.
  \item \textbf{\texttt{no\_cycle\_k\_gt\_1322}} : for \(k > 1322\),
        Hypothesis~\textsf{DerivedLargeKBound} supplies a bound
        \(n < 2^{71}\), which again contradicts the Product Bound under
        \textsf{BarinaVerification}.
\end{itemize}

% [ERIC-REVIEW] The interplay Product Bound / Baker / Barina above is a
% natural-language rendering of the proof structure inside
% ProjetCollatz.Phase59ContinuedFractions + the supporting Phase50--Phase58
% modules. Please verify that this summary does not over-simplify or
% mis-attribute the role of each hypothesis.

A detailed mathematical derivation of
\textsf{DerivedLargeKBound} from the continued-fraction expansion of
\(\log_2 3\) combined with Baker's theorem is sketched in
Section~\ref{sec:hypotheses}. Its full Lean formalisation is identified
as future work (see Section~\ref{sec:conclusion}).

\subsection{Scope and limitations}

\begin{remark}[What the theorem does \emph{not} claim]
The theorem does not assert the truth of the three hypotheses it depends
on. In particular :
\begin{itemize}
  \item \textsf{BakerSeparation} is supplied as a structural witness for
        Baker's published theorem with a weakened constant
        (see Section~\ref{sec:hypotheses}).
  \item \textsf{BarinaVerification} is supplied as a structural witness
        for Barina's published verification of the Collatz conjecture
        below \(2^{71}\).
  \item \textsf{DerivedLargeKBound} is a structural witness for a bound
        whose mathematical justification is sketched in
        Section~\ref{sec:hypotheses} but which is not, at the time of
        writing, a proven Lean theorem. Its promotion from a structure
        to a theorem is the main item of future work.
\end{itemize}
The theorem establishes a conditional non-existence : \emph{if} the three
witnesses are instantiated with valid content, \emph{then} no nontrivial
odd Collatz cycle exists.
\end{remark}
